\section{Benchmark Tasks, Data Access, and Reference Dynamics}
\label{app:benchmark}

The experiments use four modeling tasks from three system families: meal-response
modeling (T1) and absorption--action disentanglement (T2) in a glucose--insulin
system, a controlled chemical reactor, and an unfamiliar synthetic device.
A task specifies the mechanisms a model must represent; an observability tier
specifies which trajectories it may use. These are distinct from the private
simulator's state dimension or a requirement to reproduce its exact equations.

The main comparison contains nine cases: canonical and perturbed T1-easy, each
in named and obfuscated form (four cases); canonical named T1-hard, T2-easy, and
T2-hard (three cases); named CSTR-easy; and functional-description alien-device
easy. Thus, the nine cases represent three physical or synthetic families,
not nine independent systems. Only the tiers and variants used in these
experiments are described below.

\subsection{Shared data and evaluation contract}
\label{app:benchmark_contract}

Each numerical case has 16 training, four validation, and six test trajectories.
Equations and fitted parameters are shared across trajectories. Training data
support parameter estimation and proposal feedback; validation data support
model selection. Test trajectories are reserved for evaluation after models
are frozen. New initial conditions and input schedules define distinct
trajectories, rather than splitting a single trajectory into adjacent windows.

\paragraph{Targets, auxiliaries, and external inputs.}
A \emph{target} is a quantity the model must generate. During held-out free
rollout, its initial observation is available, but later target observations
cannot be used to correct the simulated state. A \emph{supplied auxiliary} is
an additional measured trajectory available throughout the prediction horizon
under the public task contract. It may enter a model's equations directly
without being assigned its own ODE. An \emph{external input} is an imposed
command or event schedule. All methods receive the same channel permissions
within a case. Latent initialization must use permitted initial information;
latent initial values are not fitted to held-out target trajectories.

Table~\ref{tab:benchmark_channels} lists the named channels. Obfuscated cases
retain the numerical data and channel roles, but replace channel names with
neutral symbols. The auxiliary contract makes the rollout \emph{conditional}
on those supplied trajectories: it does not constitute autonomous recovery of
every subsystem of the reference simulator.

\begin{table}[htbp]
\centering
\small
\setlength{\tabcolsep}{4pt}
\caption{Public channel roles in the six task--tier configurations used in the
nine-case study. $m$ denotes the meal-event channel and $u_I$ the imposed insulin
forcing. A supplied auxiliary is available throughout a rollout; a target is
predicted without later observed-value resets.}
\label{tab:benchmark_channels}
\begin{tabular}{@{}p{0.16\linewidth}p{0.17\linewidth}p{0.31\linewidth}p{0.27\linewidth}@{}}
\toprule
Task/tier & Targets & Supplied auxiliaries & External inputs \\
\midrule
T1 easy & $G_p$ & $EGP,U_{ii},E,G_t$ & $m$ \\
T1 hard & $G_p$ & $G_t$ & $m$ \\
T2 easy & $G_p,I,U$ & $EGP,U_{ii},E,G_t$ & $m,u_I$ \\
T2 hard & $G_p,I$ & $U_{ii}$ & $m,u_I$ \\
CSTR easy & $T$ & $C,T_j$ & $C_f,T_f,T_{jf}$ \\
Alien easy & $v_{01}$ & $v_{02},v_{03}$ & $u_{01}$ \\
\bottomrule
\end{tabular}
\end{table}

\paragraph{Public prompt structure.}
The prompts have six sections: task specification, available data, modeling
requirements, constraints and plausibility, intended use, and required response.
They ask for explicit analytic dynamics, a finite-dimensional causal
representation of the required mechanisms, globally shared parameters, and
causal initialization. Every modeled state needs a governing equation, every
generated process an analytic definition, and every target an observation or
governing equation. Free fitted time series, future-observation dependence,
and unrestricted neural-network vector fields are disallowed. The prompts
request coherent source/sink signs and declared state constraints, but allow
signed deviations and net rates. They do not impose a particular latent-state
count. The task-specific mechanism clauses are quoted below; the reference
equations in this appendix are explanatory and are not supplied to discovery
methods. Scientific critique and deterministic verification follow the common
methodology, rather than benchmark-specific judge prompts.

\subsection{Glucose--insulin system: T1 and T2}
\label{app:benchmark_glucose}

\paragraph{Reference system and notation.}
We generate glucose--insulin responses using a trusted implementation based on
\citet{dalla2007meal}. The reference has 12 dynamic states: plasma and tissue
glucose masses $G_p,G_t$; three gastrointestinal compartments
$Q_{\mathrm{sto1}},Q_{\mathrm{sto2}},Q_{\mathrm{gut}}$; plasma, liver, and portal
insulin quantities $I_p,I_l,I_{po}$; a secretion-response state $Y$;
two delayed insulin signals $I_1,I_d$; and a remote insulin-action state $X$.
This dimension belongs to the simulator, not to the required discovered model.
In particular, the public glucose target is mass $G_p$ in
$\mathrm{mg}\,\mathrm{kg}^{-1}$, not concentration $G=G_p/V_G$.
The insulin target is concentration $I=I_p/V_I$ in
$\mathrm{pmol}\,\mathrm{L}^{-1}$.

\paragraph{Relevant reference mechanisms.}
The glucose balances separate production, meal appearance, utilization,
excretion, and exchange:
\begin{align}
\dot G_p &= EGP + R_a - U_{ii} - E - J_{p\to t}+J_{t\to p}, \\
\dot G_t &= -U_{id}+J_{p\to t}-J_{t\to p}.
\label{eq:app_glucose_balance}
\end{align}
In the canonical system, $J_{p\to t}=k_1G_p$ and $J_{t\to p}=k_2G_t$.
Meal absorption is generated by
\begin{align}
\dot Q_{\mathrm{sto1}} &= -k_{\mathrm{gri}}Q_{\mathrm{sto1}}, \\
\dot Q_{\mathrm{sto2}} &= k_{\mathrm{gri}}Q_{\mathrm{sto1}}
-k_{\mathrm{empt}}Q_{\mathrm{sto2}}, \\
\dot Q_{\mathrm{gut}} &= k_{\mathrm{empt}}Q_{\mathrm{sto2}}
-k_{\mathrm{abs}}Q_{\mathrm{gut}},
& R_a &= \frac{f k_{\mathrm{abs}}Q_{\mathrm{gut}}}{BW}.
\end{align}
Here $k_{\mathrm{empt}}$ is the reference's stomach-content-dependent emptying
rate. Each meal of $D_j$ grams at time $t_j$ produces the reference jump
$Q_{\mathrm{sto1}}(t_j^+)=Q_{\mathrm{sto1}}(t_j^-)+1000D_j$.
The public channel \texttt{meal\_event\_g} records event amounts in grams;
it is not the unobserved appearance flux $R_a$.

Insulin-dependent disposal has dynamic memory:
\begin{align}
\dot X &= -p_{2U}X+p_{2U}(I-I_b), \\
U_{id} &= \left[\frac{(V_{m0}+V_{mx}X)G_t}{K_m+G_t}\right]_+,
& U &= U_{ii}+U_{id}, \qquad U_{ii}=F_{cns},
\end{align}
where $[a]_+=\max(a,0)$. The remote action $X$ is a deviation signal;
it is not constrained to be nonnegative merely because glucose and insulin
masses are nonnegative. Other reference contributions include
\begin{align}
EGP &= [k_{p1}-k_{p2}G_p-k_{p3}I_d-k_{p4}I_{po}]_+, \\
E &= k_{e1}[G_p-k_{e2}]_+.
\end{align}
The remaining insulin and secretion dynamics close the reference system.
For example, its plasma-insulin balance is
\begin{equation}
\dot I_p=-(m_2+m_4)I_p+m_1I_l+u_I,
\qquad I=I_p/V_I,
\end{equation}
where the imposed forcing $u_I$ has units
$\mathrm{pmol}\,\mathrm{kg}^{-1}\,\mathrm{min}^{-1}$.
These equations describe reference mechanisms, not a prescribed skeleton
or an equation-by-equation recovery target for T1 or T2.

\subsubsection{T1: meal-response modeling}
\label{app:benchmark_t1}

\paragraph{Task and use.}
T1 asks how meal timing and amount influence plasma glucose. Its intended use
is a reduced meal-response simulator: given an intake schedule and the permitted
auxiliaries, it should predict the glucose response under a new schedule.
The named public requirement is:
\begin{quote}
``a causal mechanism through which meal timing and meal amount contribute to
the observed plasma-glucose mass.''
\end{quote}
The prompt does not prescribe the three gastrointestinal compartments,
a variable named $R_a$, or any fixed number of latent states. Nor does it
require hepatic regulation, insulin-action dynamics, or a separately generated
$G_t$ state. Such mechanisms may be useful, but their absence alone is not a
violation of the T1 task specification.

\paragraph{Observability.}
T1-easy predicts $G_p$ with $EGP,U_{ii},E,G_t$ supplied. T1-hard also predicts
$G_p$, but only $G_t$ is supplied. Thus, hard means fewer auxiliaries, not
complete absence of auxiliary measurements. A candidate may choose which
auxiliaries to use, provided it explains their roles and respects the public
availability contract. Both tiers receive the meal-event channel.

\subsubsection{T2: absorption--action disentanglement}
\label{app:benchmark_t2}

\paragraph{Task and use.}
T2 extends meal-response modeling to distinguish glucose entering circulation
from glucose removal regulated by insulin. This supports interpretation of
post-meal excursions: delayed absorption and delayed insulin action can produce
different responses to a changed meal or imposed insulin input. The named
public requirements are:
\begin{quote}
``a causal pathway through which meal timing and meal amount contribute to the
observed plasma-glucose mass''; and

``a delayed insulin-action pathway through which plasma insulin regulates the
insulin-dependent contribution to total glucose disposal $U(t)$.''
\end{quote}
Unlike T1's meal clause, the second clause explicitly requires delayed action.
An instantaneous insulin-dependent coefficient alone does not supply that
memory. The task does not require the reference coordinate $X$ or its exact
functional form, and a graph path alone does not establish the correct fitted
response or a unique physical interpretation.

\paragraph{Observability.}
T2-easy must generate three targets, $G_p,I,U$, with $EGP,U_{ii},E,G_t$ supplied.
T2-hard must generate $G_p,I$, with only $U_{ii}$ supplied; total disposal is
then an internal mechanism rather than an observed target. In both tiers,
the declared external inputs are the meal-event channel and the imposed
insulin forcing. \emph{Plasma insulin $I$ is a prediction target, not a supplied
future trajectory.} The model must therefore generate the insulin signal used
by its delayed-action pathway. $U$ may be an analytic generated process rather
than an additional dynamic state. The public task does not explicitly prescribe
a meal-to-insulin or glucose-to-insulin edge; such dependencies must be proposed
from the available evidence rather than imposed as hidden reference facts.

\subsubsection{Dynamics and naming controls}
\label{app:benchmark_controls}

The T1-easy variants separate two changes. \emph{Obfuscation} replaces
physiological names with neutral channel identifiers and removes informative
units while preserving numerical trajectories, splits, and channel permissions.
It tests dependence on semantic cues; it does not by itself establish whether
an LLM memorized a model. The public task still describes the required causal
input-to-output mechanism.

The \emph{perturbed dynamics} change the two plasma--tissue exchange laws to
\begin{align}
J_{p\to t}
&= k_1G_p\left[1+0.35\tanh\!\left(
\frac{G_p-G_{p,b}}{0.20G_{p,b}}\right)\right], \\
J_{t\to p}
&= k_2G_t\left[1-0.25\tanh\!\left(
\frac{G_t-G_{t,b}}{0.20G_{t,b}}\right)\right].
\end{align}
The basal masses are $G_{p,b}$ and $G_{t,b}$. Both multipliers remain positive
and equal one at basal equilibrium, preserving exchange directions and the
basal balance while introducing state dependence. This perturbation does not
change the gastric-emptying or remote-insulin-action laws. The same exchange
flux appears with opposite signs in the two balances. The study uses the four
canonical/perturbed and named/obfuscated combinations for T1-easy; T1-hard and
T2 use the canonical named condition.

\paragraph{Trajectory schedules.}
Dalla Man trajectories last 300 minutes with one-minute sampling. Training
includes basal trajectories, single meals of 30--150\,g, delayed meals,
two-meal schedules with 60--120-minute separation, and shifts in initial glucose
or insulin. T1 adds three separated meals; T2 adds an imposed insulin pulse.
The initial $G_p$ shifts in this protocol are paired with compensating $G_t$
changes, rather than independent changes of both compartments.
Validation uses intermediate meal amounts and timings, a new two-meal schedule,
and a combined initial-state shift. Its T2 schedules contain no imposed insulin
pulse. The six test protocols include smaller/larger meals, a later meal,
a multi-meal schedule, a combined initial-state extrapolation, and an additional
task-specific schedule: four meals for T1 or a stronger insulin pulse for T2.
Post-hoc case-study interventions, such as shorter meal spacing or independently
changing $G_t(0)$, are reported separately from this fixed test suite.

\subsection{Controlled continuous stirred-tank reactor (CSTR)}
\label{app:benchmark_cstr}

\paragraph{Task and use.}
A continuously stirred-tank reactor mixes an incoming reactant stream while
a surrounding jacket exchanges heat with the reactor. Reactor temperature may
change because the feed temperature changes, an exothermic reaction releases
heat, or heat moves between the reactor and jacket. A useful model must
distinguish these contributions to support prediction under changed operating
conditions. The public task requires:
\begin{quote}
``a reactor-temperature balance that distinguishes feed transport, reaction
heat generation, and heat exchange with the jacket.''
\end{quote}
The evaluated named easy case predicts reactor temperature $T$. Reactant
concentration $C$ and jacket temperature $T_j$ are supplied auxiliaries;
feed concentration $C_f$, feed temperature $T_f$, and jacket-feed temperature
$T_{jf}$ are external inputs. Extra latent states are not compulsory when a
candidate can express the required balance using the available channels.

\paragraph{Reference dynamics.}
The three-state reference is
\begin{align}
r(C,T) &= k_0 e^{-E_R/T}C, \\
\dot C &= q(C_f-C)-r(C,T), \\
\dot T &= q(T_f-T)+\beta_r r(C,T)-h(T-T_j), \\
\dot T_j &= q_j(T_{jf}-T_j)+h_j(T-T_j).
\end{align}
Here $q,q_j$ are flow rates, $\beta_r$ converts reaction rate to heat generation,
and $h,h_j$ govern heat exchange. These positive coefficients define the
reference source/sink roles. The implemented reaction law uses
$\max(C,0)$ and $\max(T,250)$ in place of $C$ and $T$ for numerical protection.
Temperature is measured in kelvin and concentration in the supplied relative
units. The reference parameters are
$k_0=7.2\times10^{10}$, $E_R=8750$, $q=1$, $\beta_r=80$,
$h=2.5$, $q_j=1.5$, and $h_j=1.5$.

\paragraph{Trajectory schedules.}
Trajectories span 30 time units at interval 0.1. Normalized input commands
$(u_C,u_T,u_J)$ are mapped to physical inputs by
\begin{equation}
C_f=1+0.25u_C,\qquad T_f=350+10u_T,\qquad T_{jf}=330+12u_J.
\end{equation}
Training uses zero input, single-input steps and pulses, slow sinusoids,
paired input changes, and an initial-state shift. Validation uses intermediate
amplitudes and durations. The test design includes amplitude extrapolation,
faster oscillations, a delayed pulse, opposing feed/jacket temperature changes,
and initial-state extrapolation. In the easy tier, a changed physical input may
also change the supplied auxiliary trajectories; these remain part of the
conditional prediction contract.

\paragraph{Reference-data qualification.}
A diagnostic audit of a historical CSTR export found that its adaptive
integrator skipped some finite-duration forcing windows. In particular, a
350-to-345\,K feed-temperature pulse from time 10 to 15 was accompanied by an
incorrectly flat saved temperature. Reference integrations split at the forcing
boundaries instead produced a temperature decrease of approximately 4.36\,K.
Consequently, scores against affected exports describe agreement with those
stored arrays and cannot establish correct physical intervention behavior.
Any corrected release and matched reruns must be identified separately;
the task definition above does not certify the numerical integrity of a
historical export.

\subsection{Alien device: input-driven memory in an unfamiliar system}
\label{app:benchmark_alien}

\paragraph{Task and use.}
The alien device is a procedurally generated nonlinear system without a named
physical interpretation. It tests whether a model can explain a persistent
response to an external command using measured telemetry, without relying on
a familiar domain equation. The evaluated functional-description easy prompt
requires:
\begin{quote}
``input-driven memory and its causal contribution to the observed output.''
\end{quote}
The public target is $v_{01}$, the command is $u_{01}$, and $v_{02},v_{03}$ are
two supplied telemetry channels selected by a simulator-side training-design
sensitivity analysis. The prompt states the required function of memory but
withholds the reference graph, coefficients, state identities, and dimension.
It does not explicitly require the complete nonlinear-feedback architecture
of the private reference.

\paragraph{Reference dynamics.}
The selected reference has five internal coordinates $z\in\mathbb{R}^5$ and a
dynamic output $y$. Its construction has the form
\begin{align}
\dot z_i &= -d_i z_i+\sum_j A_{ij}z_j
 +\sum_{\ell}a_{i\ell}\tanh(s_{i\ell}z_{j_{i\ell}}+b_{i\ell})
 \nonumber\\
&\quad+\sum_{\ell}c_{i\ell}
 \tanh(p_{i\ell}z_{r_{i\ell}})\tanh(q_{i\ell}z_{s_{i\ell}})
 +B_i\tanh(s_u u), \\
\dot y &= -d_y y+\sum_{\ell}a^y_{\ell}\tanh(s^y_{\ell}z_{j_{\ell}})
 +\sum_{\ell}c^y_{\ell}
 \tanh(p^y_{\ell}z_{r_{\ell}})\tanh(q^y_{\ell}z_{s_{\ell}}).
\end{align}
The coefficients and interaction indices are frozen in the selected generator;
$A$ is skew-symmetric and the $d_i$ and $d_y$ are positive decay rates. The
architecture combines relaxation, internal coupling, bounded nonlinear
interactions, and command-driven forcing. Only two internal coordinates are
exposed as telemetry in the easy tier; their public names do not disclose
their private roles. Discovery is not required to reproduce all five internal
coordinates. This benchmark supplies an unfamiliar task, but neither procedural
generation nor neutral names prove that every constituent mathematical pattern
is absent from an LLM's training data.

\paragraph{Trajectory schedules.}
Trajectories span 60 time units at interval 0.1. Training includes positive and
negative pulses, steps, sinusoids, a chirp, separated pulse sequences, and
initial-state shifts. Validation uses intermediate pulse amplitudes/durations
and sinusoidal periods. Test protocols change amplitude, frequency, pulse
timing, or initial conditions beyond the training design. The model uses one
frozen parameter vector across these schedules. The relevant question is
whether its input-driven memory continues to generate the output response,
not whether it reproduces the simulator's private coordinate system.
